% Irodov, original numbered Part II problems 84--94, printed pp.72--73.
% Maxwell speed density F(v)=4/sqrt(pi) v^2/vp^3 exp[-(v/vp)^2],
% vp=sqrt(2kBT/m), exact normalization.
\thermokineticproblem{84}{72}
在麦克斯韦分布下，
求分子速率与下列特征速率相差不超过
$\delta=1.00\%$ 的分子所占比例：
(a) 最概然速率 $v_p$；
(b) 均方根速率 $v_{\rm rms}$。
\exsolution
令无量纲速率 $u=v/v_p$，
其归一化分布为
$f(u)=\frac4{\sqrt\pi}u^2e^{-u^2}$。
若速率容许偏差是特征速率 $v_*$ 的
$\delta$ 倍，则区间
$u\in[u_*(1-\delta),u_*(1+\delta)]$，
小 $\delta$ 下比例近似为
\[
P\simeq f(u_*)\,2\delta u_*
=\frac8{\sqrt\pi}\delta u_*^3e^{-u_*^2}.
\]
(a) $u_*=1$，故
\[
\boxed{P_p\simeq\frac8{e\sqrt\pi}\delta
\simeq1.66\%}.
\]
(b) 均方根速率
$v_{\rm rms}=\sqrt{3k_BT/m}$，
$u_*=\sqrt{3/2}$，
因此
\[
\boxed{P_{\rm rms}\simeq
\frac8{\sqrt\pi}\left(\frac32\right)^{3/2}
e^{-3/2}\delta\simeq1.85\%}.
\]
\excheck 用的是速率区间而非某一固定精确速率；
连续分布在单个点的概率为零，
须对小但非零宽度的区间积分。

\thermokineticproblem{85}{72}
求气体温度：
(a) 氢分子的均方根速率
比最概然速率大
$\Delta v=400\,\mathrm{m/s}$；
(b) 氧分子的速率分布函数
在 $v=420\,\mathrm{m/s}$ 处取最大值。
\exsolution
(a) 由麦克斯韦分布，
\[
v_{\rm rms}-v_p
=(\sqrt3-\sqrt2)\sqrt{\frac{k_BT}{m_H}}
=\Delta v.
\]
所以
\[
\boxed{T_a=\frac{m_H(\Delta v)^2}
{k_B(\sqrt3-\sqrt2)^2}}.
\]
用分子氢质量 $m_H\simeq2u$
可得
$\boxed{T_a\simeq380\,\mathrm K}$。
(b) 最大值出现在
$v_p=\sqrt{2k_BT/m_O}$，
故
\[
\boxed{T_b=\frac{m_Ov^2}{2k_B}
\simeq340\,\mathrm K},
\]
其中氧分子质量 $m_O\simeq32u$。
\excheck 这两问的气体摩尔质量不同；
不能把氢分子的质量用于氧分子计算。

\thermokineticproblem{86}{72}
对于氮气，求：
(a) 麦克斯韦速率密度函数
在 $v_1=300\,\mathrm{m/s}$ 和
$v_2=600\,\mathrm{m/s}$ 处数值相等的温度；
(b) 在给定速率 $v$ 处，
温度为 $T_0$ 与 $\eta T_0$ 的两个
麦克斯韦速率密度函数数值相等时的 $v$。
\exsolution
在某一温度下，
麦克斯韦速率概率密度为
\[
F_T(v)=4\pi\left(\frac m{2\pi k_BT}\right)^{3/2}
v^2e^{-mv^2/(2k_BT)}.
\]
(a) 令 $F_T(v_1)=F_T(v_2)$，
约掉温度的共同归一化系数，
得
\[
v_1^2e^{-mv_1^2/(2k_BT)}
=v_2^2e^{-mv_2^2/(2k_BT)}.
\]
取对数化简为
\[
\boxed{T=
\frac{m(v_2^2-v_1^2)}
{4k_B\ln(v_2/v_1)}}.
\]
取 $m\simeq28u$，得
$\boxed{T\simeq330\,\mathrm K}$。
(b) 比较同一个速率的两温度概率密度，
\[
T_0^{-3/2}e^{-mv^2/(2k_BT_0)}
=(\eta T_0)^{-3/2}
e^{-mv^2/(2k_B\eta T_0)}.
\]
两边取对数，得
\[
\boxed{v=
\sqrt{\frac{3k_BT_0\eta\ln\eta}
{m(\eta-1)}}}.
\]
\excheck 当 $\eta\to1$ 时，
$\eta\ln\eta/(\eta-1)\to1$，
两分布趋于相同，此式的极限只对应
由非平凡交点连续延伸到的特征速率。

\thermokineticproblem{87}{72}
在含氮分子与氧分子的混合气中，
温度为何值时氮、氧分子的最概然速率
相差 $\Delta v=30\,\mathrm{m/s}$？
\exsolution
同温混合气中，
分子量为 $m$ 的组分最概然速率为
$v_p=\sqrt{2k_BT/m}$。
氮分子比氧分子轻，
故氮的最概然速率较大。
因此
\[
\Delta v=\sqrt{2k_BT}
\left(\frac1{\sqrt{m_N}}-\frac1{\sqrt{m_O}}\right).
\]
平方求温度：
\[
\boxed{T=\frac{(\Delta v)^2}
{2k_B(1/\sqrt{m_N}-1/\sqrt{m_O})^2}}.
\]
取 $m_N\simeq28u$、$m_O\simeq32u$，
或直接以摩尔质量 $M_N=0.028\,\mathrm{kg/mol}$、
$M_O=0.032\,\mathrm{kg/mol}$ 与摩尔气体常数 $R$
计算，得
\[
T=\frac{(30\,\mathrm{m/s})^2}{
2R\left(M_N^{-1/2}-M_O^{-1/2}\right)^2}
\simeq\boxed{3.63\times10^2\,\mathrm K}.
\]
在此温度，
\[
v_{p,N}=\sqrt{\frac{2RT}{M_N}}
\simeq464.50\,\mathrm{m/s},\qquad
v_{p,O}=\sqrt{\frac{2RT}{M_O}}
\simeq434.50\,\mathrm{m/s},
\]
故 $v_{p,N}-v_{p,O}=30.00\,\mathrm{m/s}$，完成独立反代验证。
\excheck \textbf{原书答案有三个数量级的数值错误：}
底本印刷答案页第289页（PDF第298页）给出
$T=0.37\,\mathrm K$，但同时列出的符号公式
正确代入 $m_N\simeq28u$、$m_O\simeq32u$
应得约 $363\,\mathrm K$（通常可粗略写成 $0.37\,\mathrm{kK}$）。
$0.37\,\mathrm K$ 下的两种最概然速率仅相差
约 $0.96\,\mathrm{m/s}$，与题给 $30\,\mathrm{m/s}$
矛盾。已按物理方程与原始单位修正，不沿用书末错误数值。

\thermokineticproblem{88}{72}
氢气与氦气的混合物温度为
$T=300\,\mathrm K$。
在什么速率 $v$ 处，
两组分各自归一化的麦克斯韦速率分布密度
取相同值？
\exsolution
氢分子质量为 $m_H\simeq2u$，
氦原子质量为 $m_{\rm He}\simeq4u$。
两个归一化速率概率密度相等意味着
\[
m_H^{3/2}v^2e^{-m_Hv^2/(2k_BT)}
=m_{\rm He}^{3/2}v^2e^{-m_{\rm He}v^2/(2k_BT)}.
\]
除去共同的 $v^2$（取 $v>0$），
两边取自然对数：
\[
\frac{(m_{\rm He}-m_H)v^2}{2k_BT}
=\frac32\ln\frac{m_{\rm He}}{m_H}.
\]
所以
\[
\boxed{v=
\sqrt{\frac{3k_BT\ln(m_{\rm He}/m_H)}
{m_{\rm He}-m_H}}}.
\]
代入后
$\boxed{v\simeq1.6\,\mathrm{km/s}}$。
\excheck 这里比较的是\emph{各自归一化的函数}
$F(v)$，不包含混合气未知的组分浓度权重。

\thermokineticproblem{89}{72}
分子质量为 $m$ 的理想气体中，
对于固定速率区间 $[v,v+\mathrm dv]$，
在何种温度 $T$ 下，
速率处于该区间的分子数量最多？
\exsolution
归一化麦克斯韦速率密度
\[
F_T(v)=
4\pi\left(\frac m{2\pi k_BT}\right)^{3/2}
v^2e^{-mv^2/(2k_BT)}.
\]
当 $v$、$\mathrm dv$ 固定时，
需最大化温度函数
$T^{-3/2}\exp[-mv^2/(2k_BT)]$。
取对数并求导：
\[
\frac{\mathrm d}{\mathrm dT}\ln F_T(v)
=-\frac3{2T}+\frac{mv^2}{2k_BT^2}=0.
\]
得到唯一极大值
\[
\boxed{T=\frac{mv^2}{3k_B}}.
\]
\excheck 这是固定某一速率窗口对温度求最大；
不同于固定温度下对速率求最概然值
$v_p=\sqrt{2k_BT/m}$。

\thermokineticproblem{90}{73}
求气体中分子速度满足以下条件的概率：
$x$ 方向速度分量属于
$[v_x,v_x+\mathrm dv_x]$，
同时垂直 $x$ 方向的速度分量模
$v_\perp=\sqrt{v_y^2+v_z^2}$
属于 $[v_\perp,v_\perp+\mathrm dv_\perp]$。
分子质量 $m$，温度 $T$。
\exsolution
三维麦克斯韦速度密度为
\[
f(\bm v)=
\left(\frac m{2\pi k_BT}\right)^{3/2}
e^{-m(v_x^2+v_y^2+v_z^2)/(2k_BT)}.
\]
在垂直的 $(v_y,v_z)$ 平面上引入极坐标，
$\mathrm dv_y\mathrm dv_z
=v_\perp\,\mathrm dv_\perp\,\mathrm d\varphi$。
把极角 $\varphi\in[0,2\pi)$ 积分掉：
\[
\boxed{\frac{\mathrm dN}{N}
=2\pi\left(\frac m{2\pi k_BT}\right)^{3/2}
v_\perp e^{-m(v_x^2+v_\perp^2)/(2k_BT)}
\,\mathrm dv_x\,\mathrm dv_\perp}.
\]
\excheck 变量 $v_\perp$ 是二维分量\emph{模}，
须包含极坐标雅可比因子 $2\pi v_\perp$。

\thermokineticproblem{91}{73}
从麦克斯韦分布直接计算
理想气体分子速度 $x$ 分量的平均值
$\overline{v_x}$，
以及其绝对值的平均值
$\overline{|v_x|}$。
\exsolution
一维速度分布为
\[
f_x(v_x)=\sqrt{\frac m{2\pi k_BT}}
e^{-mv_x^2/(2k_BT)}.
\]
因为该函数为偶函数，
$v_xf_x(v_x)$ 为奇函数，
所以
\[
\boxed{\overline{v_x}
=\int_{-\infty}^{\infty}v_xf_x(v_x)\,\mathrm dv_x=0}.
\]
绝对值平均为
\[
\overline{|v_x|}
=2\sqrt{\frac m{2\pi k_BT}}
\int_0^\infty
v_xe^{-mv_x^2/(2k_BT)}\,\mathrm dv_x.
\]
令
$u=mv_x^2/(2k_BT)$，
$v_x\mathrm dv_x=(k_BT/m)\mathrm du$，
得
\[
\boxed{\overline{|v_x|}
=\sqrt{\frac{2k_BT}{\pi m}}}.
\]
\excheck 平均的有符号分量为零，
但其绝对值平均严格为正；
不能将两种平均混淆。

\thermokineticproblem{92}{73}
利用麦克斯韦分布求
气体分子 $x$ 方向速度分量平方的平均值
$\overline{v_x^2}$，气体温度 $T$、分子质量 $m$。
\exsolution
代入归一化一维高斯分布：
\[
\overline{v_x^2}
=\sqrt{\frac m{2\pi k_BT}}
\int_{-\infty}^{\infty}
v_x^2e^{-mv_x^2/(2k_BT)}\,\mathrm dv_x.
\]
利用高斯积分恒等式
\[
\int_{-\infty}^{\infty}x^2e^{-ax^2}\,\mathrm dx
=\frac{\sqrt\pi}{2a^{3/2}},\qquad a>0,
\]
取 $a=m/(2k_BT)$，化简为
\[
\boxed{\overline{v_x^2}=\frac{k_BT}{m}}.
\]
\excheck 三个正交速度分量独立且等方差，
所以 $\overline{v^2}=3k_BT/m$，与均分定理一致。

\thermokineticproblem{93}{73}
在温度为 $T$、数密度为 $n$、
分子质量为 $m$ 的理想气体中，
根据麦克斯韦分布求
单位墙面积每单位时间受到分子撞击的次数
$\Phi$。
\exsolution
取墙法线为 $x$ 轴，
只有朝墙运动、满足 $v_x>0$ 的分子
能撞击墙。
单位面积每秒碰撞数是正向粒子通量
\[
\Phi=n\int_0^\infty v_x f_x(v_x)\,\mathrm dv_x,
\quad
f_x(v_x)=\sqrt{\frac m{2\pi k_BT}}
e^{-mv_x^2/(2k_BT)}.
\]
计算积分：
\[
\int_0^\infty v_x f_x(v_x)\,\mathrm dv_x
=\sqrt{\frac{k_BT}{2\pi m}},
\]
所以
\[
\boxed{\Phi=n\sqrt{\frac{k_BT}{2\pi m}}
=\frac14n\sqrt{\frac{8k_BT}{\pi m}}
=\frac14n\bar v}.
\]
\excheck 因为正、负法向速度各占一半，
而撞击通量还按速率加权，
不能只用 $n\bar v/2$。

\thermokineticproblem{94}{73}
从麦克斯韦分布直接求
数密度为 $n$、温度为 $T$ 的气体
对静止平壁产生的压强 $p$，
单分子质量为 $m$。
\exsolution
每个以正法向分量 $v_x>0$ 撞墙的
分子在弹性反射中传给墙
法向冲量 $2mv_x$；
相应撞击通量微元为
$n v_x f_x(v_x)\,\mathrm dv_x$。
所以
\[
p=2mn\int_0^\infty v_x^2 f_x(v_x)\,\mathrm dv_x.
\]
由于高斯分布为偶函数，
\[
2\int_0^\infty v_x^2 f_x(v_x)\,\mathrm dv_x
=\int_{-\infty}^\infty v_x^2 f_x(v_x)\,\mathrm dv_x
=\overline{v_x^2}.
\]
上一题得到
$\overline{v_x^2}=k_BT/m$，
因此
\[
\boxed{p=mn\overline{v_x^2}=nk_BT}.
\]
\excheck 这是理想气体状态方程的动理论形式；
压强取决于法向动量通量，而不是
墙上单次碰撞冲量大小。
